\documentclass{article}
\usepackage[T1]{fontenc}
\usepackage{lmodern}
\usepackage{graphicx}
\usepackage{amsmath,amssymb}
\usepackage[a4paper,margin=2cm]{geometry}
\usepackage[absolute,overlay]{textpos}
\setlength{\TPHorizModule}{1mm}
\setlength{\TPVertModule}{1mm}
\usepackage[indonesian]{babel}
\usepackage{hyperref}
\setlength{\emergencystretch}{2em}
% unit: Halaman 4

\begin{document}

% === Halaman 4 (python/native) ===

4

• Kelemahan cipher abjad-tunggal: tidak dapat menyembunyikan 

hubungan statistik antara plainteks dengan cipherteks.

   -  Huruf yang sama dienkripsi menjadi huruf cipherteks yang sama

   -  Huruf yang sering muncul di dalam plainteks, juga sering muncul 

di dalam huruf cipherteksnya yang berkoresponden.

    - artinya struktur statistik bahasa masih muncul di dalam cipherteks

• Oleh karena itu, cipherteks dapat didekripsi tanpa mengetahui 

kuncinya sekalipun dengan menggunakan teknik analisis frekuensi.

Kriptanalisis Cipher Abjad-Tunggal

\end{document}
